\begin{proof}[Proof of \cref{obj:thm:weighted-packet-construction}]
Fix \(\kappa\in[0,2]\). We use the filter, taper parameter, kernels, and packet from \cref{obj:def:packet-handle}. All constants below are independent of the integer \(m\ge4\).

\begin{enumerate}
\item \emph{Filter properties and cited kernel estimates.}
The standard smooth function
\[
\vartheta(v)=
\begin{cases}
e^{-1/v},&v>0,\\
0,&v\le0
\end{cases}
\qquad(v\in\mathbb R)
\]
gives
\[
\eta(s)=\vartheta\bigl((s-9/8)(15/8-s)\bigr).
\]
Consequently, \(\eta\) is smooth, \(0\le\eta\le1\), and
\[
\eta(s)\ne0\ \Longleftrightarrow\ 9/8<s<15/8.
\]
In particular,
\[
\eta(j/m)\ne0\quad\Longrightarrow\quad m<j<2m.
\]
On the central spectral interval,
\[
(s-9/8)(15/8-s)\ge\frac1{64}
\quad(5/4\le s\le7/4),
\]
so
\[
\eta(s)\ge c_\eta,\qquad c_\eta:=e^{-64}>0.
\]

For Jacobi parameters \(a_{\mathrm J},b_{\mathrm J}>-1/2\), define
\[
\mathcal K_m^{a_{\mathrm J},b_{\mathrm J}}(z_1,z_2)
:=
\sum_{j=0}^{2m}\eta(j/m)
\frac{\mathsf J_j^{(a_{\mathrm J},b_{\mathrm J})}(z_1)
      \mathsf J_j^{(a_{\mathrm J},b_{\mathrm J})}(z_2)}
     {H_j^{a_{\mathrm J},b_{\mathrm J}}}
\qquad(z_1,z_2\in\mathbb R),
\]
and, for \(z,z_1,z_2\in[-1,1]\), define
\[
\mathcal W_{a_{\mathrm J},b_{\mathrm J}}(m;z)
:=
(1-z+m^{-2})^{a_{\mathrm J}+1/2}
(1+z+m^{-2})^{b_{\mathrm J}+1/2},
\qquad
\delta_{\mathrm{ang}}(z_1,z_2)
:=
|\arccos z_1-\arccos z_2|.
\]
The spectral support permits truncation at degree \(2m\). Petrushev--Xu use squared norms normalized by the total Jacobi weight mass
\citep[pp.~1--2, equation~(1.2)]{PetrushevXu2005}. In the present notation, their norms and filtered kernel are
\[
\mathsf h_{\mathrm{PX},j}^{a_{\mathrm J},b_{\mathrm J}}
:=\frac{H_j^{a_{\mathrm J},b_{\mathrm J}}}
        {H_0^{a_{\mathrm J},b_{\mathrm J}}}
\qquad(j\ge0),
\]
\[
\begin{split}
\mathcal K_{\mathrm{PX},m}^{a_{\mathrm J},b_{\mathrm J}}(z_1,z_2)
&:=\sum_{j=0}^{2m}\eta(j/m)
\frac{\mathsf J_j^{(a_{\mathrm J},b_{\mathrm J})}(z_1)
      \mathsf J_j^{(a_{\mathrm J},b_{\mathrm J})}(z_2)}
     {\mathsf h_{\mathrm{PX},j}^{a_{\mathrm J},b_{\mathrm J}}}\\
&=H_0^{a_{\mathrm J},b_{\mathrm J}}
  \mathcal K_m^{a_{\mathrm J},b_{\mathrm J}}(z_1,z_2)
\qquad(z_1,z_2\in\mathbb R).
\end{split}
\]
Since \(H_0^{a_{\mathrm J},b_{\mathrm J}}>0\) depends only on the fixed Jacobi parameters, dividing their localization estimate by this factor absorbs its reciprocal into the localization constant below. The localization estimate of
\citep[Theorem~2.4, equation~(2.14)]{PetrushevXu2005}
provides, for every \(\tau>0\), a positive constant
\(C_{\mathrm{loc},\tau}^{a_{\mathrm J},b_{\mathrm J}}\) such that
\[
|\mathcal K_m^{a_{\mathrm J},b_{\mathrm J}}(z_1,z_2)|
\le
\frac{C_{\mathrm{loc},\tau}^{a_{\mathrm J},b_{\mathrm J}}m}
{\sqrt{\mathcal W_{a_{\mathrm J},b_{\mathrm J}}(m;z_1)
       \mathcal W_{a_{\mathrm J},b_{\mathrm J}}(m;z_2)}
 (1+m\delta_{\mathrm{ang}}(z_1,z_2))^\tau}.
\]
The local angular estimate of
\citep[Theorem~2.2, equation~(2.4)]{https://people.math.sc.edu/pencho/Publications/kpx-06-28-06-web.pdf},
with its fixed neighborhood constant equal to \(1\), provides a positive
\(C_{\mathrm{lip},\tau}^{a_{\mathrm J},b_{\mathrm J}}\) such that
\[
\begin{split}
&|\mathcal K_m^{a_{\mathrm J},b_{\mathrm J}}(z_1,z_2)
 -\mathcal K_m^{a_{\mathrm J},b_{\mathrm J}}(\xi,z_2)|\\
&\qquad\le
\frac{C_{\mathrm{lip},\tau}^{a_{\mathrm J},b_{\mathrm J}}m^2
      \delta_{\mathrm{ang}}(z_1,\xi)}
{\sqrt{\mathcal W_{a_{\mathrm J},b_{\mathrm J}}(m;z_2)
       \mathcal W_{a_{\mathrm J},b_{\mathrm J}}(m;\zeta)}
 (1+m\delta_{\mathrm{ang}}(z_2,\zeta))^\tau},
\end{split}
\]
whenever \(z_1,z_2,\xi,\zeta\in[-1,1]\) and
\[
\delta_{\mathrm{ang}}(z_1,\xi)\le m^{-1},
\qquad
\delta_{\mathrm{ang}}(\zeta,\xi)\le m^{-1}.
\]
The parameter and filter hypotheses hold for both Jacobi systems used below.
% lean: CausalSmith.Stat.NoisydoseWeakdesignTransition.FilteredJacobiLipschitz

\item \emph{Norms and uniform gamma-ratio bounds.}
Jacobi orthogonality and the positive squared-norm formula are supplied by
\citep[Table~18.3.1, Jacobi row and caption]{https://dlmf.nist.gov/18.3}:
\[
H_j^{a_{\mathrm J},b_{\mathrm J}}
=
\frac{2^{a_{\mathrm J}+b_{\mathrm J}+1}}
     {2j+a_{\mathrm J}+b_{\mathrm J}+1}
\frac{\Gamma(j+a_{\mathrm J}+1)\Gamma(j+b_{\mathrm J}+1)}
     {\Gamma(j+1)\Gamma(j+a_{\mathrm J}+b_{\mathrm J}+1)}
>0.
\]
Orthogonality applies against every polynomial of degree less than \(j\).

For fixed positive shifts \(\lambda_1,\lambda_2\), the asymptotic in
\citep[Equation~5.11.12]{https://dlmf.nist.gov/5.11.E12}
gives an integer cutoff \(j_\Gamma\) such that
\[
\frac12
\le
\frac{\Gamma(j+\lambda_1)}
     {\Gamma(j+\lambda_2)j^{\lambda_1-\lambda_2}}
\le2
\qquad(j\ge j_\Gamma,\ j\ge1).
\]
All these normalized ratios are positive. Thus the constants
\[
\begin{split}
c_\Gamma(\lambda_1,\lambda_2)
&:=
\min\left(
\{1/2\}\cup
\left\{
\frac{\Gamma(j+\lambda_1)}
     {\Gamma(j+\lambda_2)j^{\lambda_1-\lambda_2}}
:1\le j<j_\Gamma
\right\}\right),\\
C_\Gamma(\lambda_1,\lambda_2)
&:=
\max\left(
\{2\}\cup
\left\{
\frac{\Gamma(j+\lambda_1)}
     {\Gamma(j+\lambda_2)j^{\lambda_1-\lambda_2}}
:1\le j<j_\Gamma
\right\}\right)
\end{split}
\]
are positive and satisfy
\[
c_\Gamma(\lambda_1,\lambda_2)j^{\lambda_1-\lambda_2}
\le
\frac{\Gamma(j+\lambda_1)}{\Gamma(j+\lambda_2)}
\le
C_\Gamma(\lambda_1,\lambda_2)j^{\lambda_1-\lambda_2}
\qquad(j\ge1).
\]
The singleton sets in these definitions also cover an empty initial segment.
% lean: CausalSmith.Stat.NoisydoseWeakdesignTransition.gamma_ratio_integer_power_bounds

\item \emph{Endpoint normalization for \(\kappa>0\).}
Suppose first that \(\kappa>0\), and set
\[
b:=\frac{\kappa-1}{2}>-\frac12.
\]
The standard endpoint evaluation, followed by the gamma recurrence, gives
\[
\mathsf J_j^{(4,b)}(-1)
=
\frac{(-1)^j}{j!}\prod_{i=0}^{j-1}(b+1+i)
=
(-1)^j\frac{\Gamma(j+b+1)}
{\Gamma(b+1)\Gamma(j+1)}
\qquad(j\ge0),
\]
where the empty product is \(1\). Every factor in the product is positive, so the endpoint value is nonzero. Combining this identity with the norm formula yields
\[
\frac{\mathsf J_j^{(4,b)}(-1)^2}{H_j^{4,b}}
=
\frac{2j+b+5}{2^{b+5}\Gamma(b+1)^2}
\frac{\Gamma(j+b+1)}{\Gamma(j+1)}
\frac{\Gamma(j+b+5)}{\Gamma(j+5)}.
\]
For \(j\ge1\),
\[
2j\le2j+b+5\le(b+7)j.
\]
Applying the preceding gamma bounds to the shift pairs
\((b+1,1)\) and \((b+5,5)\), define
\[
\begin{split}
c_{\mathrm{spec,end}}
&:=
\frac{2c_\Gamma(b+1,1)c_\Gamma(b+5,5)}
     {2^{b+5}\Gamma(b+1)^2},\\
C_{\mathrm{spec,end}}
&:=
\frac{(b+7)C_\Gamma(b+1,1)C_\Gamma(b+5,5)}
     {2^{b+5}\Gamma(b+1)^2}.
\end{split}
\]
These constants are positive, and \(2b+1=\kappa\) gives
\[
c_{\mathrm{spec,end}}j^\kappa
\le
\frac{\mathsf J_j^{(4,b)}(-1)^2}{H_j^{4,b}}
\le
C_{\mathrm{spec,end}}j^\kappa
\qquad(j\ge1).
\]

Every summand of \(K_m(-1)\) is nonnegative. Strict positivity follows already from the degree
\[
\nu_{\mathrm{pos}}(m):=m+\lfloor m/2\rfloor,
\qquad
\frac98<\frac{\nu_{\mathrm{pos}}(m)}m<\frac{15}8
\quad(m\ge4),
\]
whose filter value and squared endpoint coefficient are positive.

For the quantitative lower bound, use the distinct degrees
\[
\nu_{\mathrm{end}}(m,i)
:=
m+\lfloor m/4\rfloor+1+i,
\qquad
0\le i<\lfloor m/4\rfloor.
\]
They satisfy
\[
\frac54\le\frac{\nu_{\mathrm{end}}(m,i)}m\le\frac74,
\qquad
m\le\nu_{\mathrm{end}}(m,i)\le2m.
\]
Moreover,
\[
m\le4\lfloor m/4\rfloor+3
\le7\lfloor m/4\rfloor
\qquad(m\ge4).
\]
Each selected summand is therefore at least
\(c_\eta c_{\mathrm{spec,end}}m^\kappa\). For the upper bound, inactive summands vanish, whereas each active summand is at most
\(C_{\mathrm{spec,end}}2^\kappa m^\kappa\); there are at most \(2m+1\le3m\) summands. Hence, with
\[
c_{\mathrm{end}}
:=\frac{c_\eta c_{\mathrm{spec,end}}}{7}>0,
\qquad
C_{\mathrm{end}}
:=3\,2^\kappa C_{\mathrm{spec,end}}>0,
\]
we obtain
\[
0<c_{\mathrm{end}}m^{\kappa+1}
\le K_m(-1)\le C_{\mathrm{end}}m^{\kappa+1}
\qquad(m\ge4).
\]
% lean: CausalSmith.Stat.NoisydoseWeakdesignTransition.packetKernel_endpoint_power_bounds

\item \emph{Endpoint localization.}
Define the polynomial coordinate and taper base on the whole line by
\[
z_{\mathrm{end}}(u):=2u^2-1,
\qquad
\Delta(u):=1-u^2
\qquad(u\in\mathbb R).
\]
For \(|u|\le1\), define
\[
\mathcal B_{\mathrm{end}}(m,u)
:=
\sqrt{\mathcal W_{4,b}(m;-1)
      \mathcal W_{4,b}(m;z_{\mathrm{end}}(u))}.
\]
Here \(0\le\Delta(u)\le1\), and
\[
\begin{split}
\mathcal W_{4,b}(m;-1)
&=(2+m^{-2})^{9/2}m^{-\kappa}
\ge m^{-\kappa},\\
\mathcal W_{4,b}(m;z_{\mathrm{end}}(u))
&=(2\Delta(u)+m^{-2})^{9/2}
  (2u^2+m^{-2})^{\kappa/2}.
\end{split}
\]
The inequalities
\[
\Delta(u)^8
\le\Delta(u)^{9/2}
\le(2\Delta(u)+m^{-2})^{9/2},
\qquad
m^{-\kappa}\le(2u^2+m^{-2})^{\kappa/2}
\]
therefore imply
\[
\mathcal B_{\mathrm{end}}(m,u)
\ge\Delta(u)^4m^{-\kappa}.
\]
Also,
\[
\delta_{\mathrm{ang}}(-1,z_{\mathrm{end}}(u))
=2\arcsin|u|
\ge|u|.
\]
Indeed, the equality follows from
\(\arccos(2u^2-1)=\pi-2\arcsin|u|\), and the inequality follows from
\(\sin v\le v\) for \(v\ge0\).

Use localization with \(\tau=\kappa+2\). On \([-1,1]\), positivity of the normalizer gives
\[
\begin{split}
|\psi_m(u)|
&=
\frac{\Delta(u)^4|K_m(z_{\mathrm{end}}(u))|}{K_m(-1)}\\
&\le
\frac{\Delta(u)^4 C_{\mathrm{loc},\kappa+2}^{4,b}m}
{\mathcal B_{\mathrm{end}}(m,u)
 (1+m\delta_{\mathrm{ang}}(-1,z_{\mathrm{end}}(u)))^{\kappa+2}
 K_m(-1)}\\
&\le
C_{\mathrm{dec,end}}(1+m|u|)^{-(\kappa+2)},
\end{split}
\qquad
C_{\mathrm{dec,end}}
:=
\frac{C_{\mathrm{loc},\kappa+2}^{4,b}}{c_{\mathrm{end}}}>0.
\]
Outside this interval \(\psi_m=0\), so the same bound holds for every real \(u\).
% lean: CausalSmith.Stat.NoisydoseWeakdesignTransition.packetPsi_endpoint_localization

\item \emph{Endpoint derivative and continuous differentiability.}
Fix \(0<|u|<1\). In the local angular estimate take the fixed second argument to be \(-1\), the reference and auxiliary arguments both to be \(z_{\mathrm{end}}(u)\), and the varying argument to be \(z_{\mathrm{end}}(v)\). For \(v\) sufficiently close to \(u\), all arguments belong to \([-1,1]\) and the required angular distance is at most \(m^{-1}\). Symmetry of the finite kernel sum then gives
\[
\begin{split}
&|K_m(z_{\mathrm{end}}(v))-K_m(z_{\mathrm{end}}(u))|\\
&\quad\le
\frac{C_{\mathrm{lip},1}^{4,b}m^2}
{\mathcal B_{\mathrm{end}}(m,u)
 (1+m\delta_{\mathrm{ang}}(-1,z_{\mathrm{end}}(u)))}
|\arccos z_{\mathrm{end}}(v)-\arccos z_{\mathrm{end}}(u)|.
\end{split}
\]
Dividing by \(|v-u|\) and passing to the limit is valid because the kernel is a finite polynomial sum. The identities
\[
1-z_{\mathrm{end}}(u)^2=4u^2\Delta(u),
\qquad
\left|\frac{d}{du}\arccos z_{\mathrm{end}}(u)\right|
=\frac2{\sqrt{\Delta(u)}}
\]
yield
\[
|(K_m\circ z_{\mathrm{end}})'(u)|
\le
\frac{2C_{\mathrm{lip},1}^{4,b}m^2}
{\mathcal B_{\mathrm{end}}(m,u)
 (1+m\delta_{\mathrm{ang}}(-1,z_{\mathrm{end}}(u)))
 \sqrt{\Delta(u)}}
\le
\frac{2C_{\mathrm{lip},1}^{4,b}m^2}
{\mathcal B_{\mathrm{end}}(m,u)\sqrt{\Delta(u)}}.
\]
Localization with exponent \(1\) likewise gives
\[
|K_m(z_{\mathrm{end}}(u))|
\le
\frac{C_{\mathrm{loc},1}^{4,b}m}
{\mathcal B_{\mathrm{end}}(m,u)}.
\]

For the derivative assembly, use the cubed taper. Since
\[
\Delta(u)^6
\le\Delta(u)^{9/2}
\le(2\Delta(u)+m^{-2})^{9/2},
\]
the same weight comparison as before gives
\[
\frac{\Delta(u)^3}{\mathcal B_{\mathrm{end}}(m,u)}
\le m^\kappa.
\]
Furthermore,
\[
\frac{\Delta(u)^4}{\sqrt{\Delta(u)}}\le\Delta(u)^3,
\]
because \(0<\Delta(u)\le1\). The product rule gives
\[
\psi_m'(u)=
\frac{-8u\Delta(u)^3K_m(z_{\mathrm{end}}(u))
+\Delta(u)^4(K_m\circ z_{\mathrm{end}})'(u)}
{K_m(-1)}.
\]
Using \(|u|\le1\), the preceding inequalities, and the normalization lower bound,
\[
\begin{split}
|\psi_m'(u)|
&\le
\frac{
8C_{\mathrm{loc},1}^{4,b}m
 \Delta(u)^3/\mathcal B_{\mathrm{end}}(m,u)
+
2C_{\mathrm{lip},1}^{4,b}m^2
 \Delta(u)^4/(\mathcal B_{\mathrm{end}}(m,u)\sqrt{\Delta(u)})
}{K_m(-1)}\\
&\le
\frac{8C_{\mathrm{loc},1}^{4,b}m^{\kappa+1}
      +2C_{\mathrm{lip},1}^{4,b}m^{\kappa+2}}
     {c_{\mathrm{end}}m^{\kappa+1}}\\
&\le C_{\mathrm{der,end}}m,
\end{split}
\qquad
C_{\mathrm{der,end}}
:=
\frac{8C_{\mathrm{loc},1}^{4,b}
      +2C_{\mathrm{lip},1}^{4,b}}{c_{\mathrm{end}}}>0.
\]

To obtain continuous differentiability on the whole line, define
\[
\mathcal T(u):=
\begin{cases}
(1-u^2)^4,&|u|\le1,\\
0,&|u|>1.
\end{cases}
\]
For every \(u\in\mathbb R\),
\[
\mathcal T(u)
=
\frac{(1-u^2)^4+(1-u^2)|1-u^2|^3}{2}.
\]
The function \(v\mapsto|v|^3\) is continuously differentiable, with derivative \(3v|v|\). Thus \(\mathcal T\) is continuously differentiable, and
\[
\psi_m(u)
=
\mathcal T(u)\frac{K_m(z_{\mathrm{end}}(u))}{K_m(-1)}
\qquad(u\in\mathbb R)
\]
is a product of continuously differentiable functions. Outside \([-1,1]\) the packet is locally zero, so
\[
\psi_m'(u)=0\qquad(|u|>1).
\]
Continuity of \(\psi_m'\) extends the derivative bound from the complement of
\(\{-1,0,1\}\) to all of \(\mathbb R\).
% lean: CausalSmith.Stat.NoisydoseWeakdesignTransition.packetPsi_endpoint_derivative_bound

\item \emph{Endpoint mass bounds and common constant.}
Set
\[
C_{\mathrm{env,end}}
:=
\max\{C_{\mathrm{dec,end}},C_{\mathrm{der,end}}\},
\qquad
C_{\mathrm{mid,end}}:=2C_{\mathrm{env,end}}.
\]
Both constants are positive. Since the coordinate \(z_{\mathrm{end}}\), taper, and support are invariant under reflection,
\[
\psi_m(-u)=\psi_m(u)\qquad(u\in\mathbb R).
\]
The envelope and \(1+m|u|\ge1\) give
\[
|\psi_m(u)|\le C_{\mathrm{env,end}}\le C_{\mathrm{mid,end}},
\qquad
|\psi_m'(u)|\le C_{\mathrm{mid,end}}m.
\]

For \(0\le u\le1\), the inequality
\((mu)^\kappa\le(1+mu)^\kappa\) gives
\[
u^\kappa(1+mu)^{-(\kappa+2)}
\le m^{-\kappa}(1+mu)^{-2}.
\]
The quadratic envelope has the explicit integral
\[
\int_0^1(1+mu)^{-2}\,du
=
\left[-\frac1{m(1+mu)}\right]_0^1
=
\frac1m-\frac1{m(1+m)}
\le\frac1m.
\]
By evenness,
\[
\begin{split}
\int_{-1}^1|u|^\kappa|\psi_m(u)|\,du
&=
2\int_0^1u^\kappa|\psi_m(u)|\,du\\
&\le
2C_{\mathrm{env,end}}m^{-\kappa}
\int_0^1(1+mu)^{-2}\,du\\
&\le C_{\mathrm{mid,end}}m^{-\kappa-1}.
\end{split}
\]
Continuity ensures integrability of these functions on the compact interval. Pointwise,
\[
\psi_m(u)^2
=
|\psi_m(u)|^2
\le C_{\mathrm{mid,end}}|\psi_m(u)|,
\]
and hence
\[
\int_{-1}^1|u|^\kappa\psi_m(u)^2\,du
\le
C_{\mathrm{mid,end}}
\int_{-1}^1|u|^\kappa|\psi_m(u)|\,du
\le C_{\mathrm{mid,end}}^2m^{-\kappa-1}.
\]
Define
\[
C_{\mathrm{fin,end}}
:=
\max\{C_{\mathrm{mid,end}},C_{\mathrm{mid,end}}^2\}>0.
\]
The supremum, derivative, absolute-mass, and squared-mass bounds all hold with this common constant.
% lean: CausalSmith.Stat.NoisydoseWeakdesignTransition.packetendpoint_bounds

\item \emph{Endpoint value, support, and moments.}
The definition and positive normalizer give
\[
\psi_m(u)=0\quad(|u|>1),
\qquad
\psi_m(0)=\frac{K_m(-1)}{K_m(-1)}=1.
\]

Let \(s\) be an integer with \(0\le s\le m\). Every active spectral degree exceeds \(m\), while the polynomial
\(((1+z)/2)^s\) has degree \(s\). Termwise Jacobi orthogonality therefore gives
\[
\int_{-1}^1
(1-z)^4(1+z)^bK_m(z)
\left(\frac{1+z}{2}\right)^s\,dz=0.
\]
For \(0<u<1\), the substitution \(z=2u^2-1\) has Jacobian \(4u\), and
\[
\begin{split}
&(1-(2u^2-1))^4(1+(2u^2-1))^b
\left(\frac{1+(2u^2-1)}2\right)^s(4u)\\
&\qquad=
2^{6+b}u^\kappa u^{2s}(1-u^2)^4.
\end{split}
\]
Consequently,
\[
\begin{split}
\int_0^1u^\kappa u^{2s}\psi_m(u)\,du
&=
\frac1{2^{6+b}K_m(-1)}
\int_{-1}^1
(1-z)^4(1+z)^bK_m(z)
\left(\frac{1+z}{2}\right)^s\,dz\\
&=0.
\end{split}
\]
The integrand of the full even moment is even, so reflection gives
\[
\int_{-1}^1|u|^\kappa u^{2s}\psi_m(u)\,du
=
2\int_0^1u^\kappa u^{2s}\psi_m(u)\,du=0.
\]
For an odd integer \(j\), reflection changes the sign of the integrand and gives
\[
\int_{-1}^1|u|^\kappa u^j\psi_m(u)\,du
=
-\int_{-1}^1|u|^\kappa u^j\psi_m(u)\,du
=0.
\]
If \(0\le j<m\) is even, writing \(j=2s\) gives \(s\le m\), so its moment vanishes as well. Thus the positive-exponent branch satisfies all assertions with
\[
C:=C_{\mathrm{fin,end}},
\qquad c:=1.
\]
% lean: CausalSmith.Stat.NoisydoseWeakdesignTransition.packetendpoint_bounds

\item \emph{Interior spectral lower bound.}
In the remaining case, \(0\le\kappa\) and \(\kappa\le0\), hence \(\kappa=0\). We use the symmetric Jacobi system with parameters \((4,4)\).

First, the norm formula and gamma bounds for the shift pairs \((5,1)\) and \((5,9)\) give
\[
\begin{split}
H_j^{4,4}
&=
\frac{512}{2j+9}
\frac{\Gamma(j+5)}{\Gamma(j+1)}
\frac{\Gamma(j+5)}{\Gamma(j+9)}\\
&\le
\frac{512}{2j+9}
C_\Gamma(5,1)j^4 C_\Gamma(5,9)j^{-4}
\le\frac{C_{\mathrm{norm,int}}}{j}
\qquad(j\ge1),
\end{split}
\]
where
\[
C_{\mathrm{norm,int}}
:=512C_\Gamma(5,1)C_\Gamma(5,9)>0.
\]

The binomial representation in \cref{obj:lem:jacobi-binomial}, evaluated at zero, gives
\[
\mathsf J_{2s}^{(4,4)}(0)
=
4^{-s}
\sum_{i=0}^{2s}
\binom{2s+4}{i}\binom{2s+4}{2s-i}(-1)^{2s-i}.
\]
The sum is the coefficient of degree \(2s\) in
\((1+z)^{2s+4}(1-z)^{2s+4}=(1-z^2)^{2s+4}\). Therefore
\[
\mathsf J_{2s}^{(4,4)}(0)
=
(-1)^s4^{-s}\binom{2s+4}{s}
\qquad(s\ge0).
\]
In particular, these even-degree center values are nonzero.

Wallis' inequality, in its finite-product form, states
\[
\prod_{i=1}^s\frac{(2i)^2}{(2i-1)(2i+1)}
=
\frac{2^{4s}(s!)^4}{((2s)!)^2(2s+1)}
\le\frac\pi2
\qquad(s\ge0).
\]
The factorial identity
\[
\frac{2^{4s}(s!)^4}{((2s)!)^2(2s+1)}
\left(4^{-s}\binom{2s}{s}\right)^2(2s+1)=1
\]
then gives, for \(s\ge1\),
\[
\left(4^{-s}\binom{2s}{s}\right)^2
\ge\frac2{\pi(2s+1)}
\ge\frac2{3\pi s}.
\]
Since \(\binom{2s+4}{s}\ge\binom{2s}{s}\),
\[
\mathsf J_{2s}^{(4,4)}(0)^2\ge\frac2{3\pi s}.
\]
Together with
\[
0<H_{2s}^{4,4}
\le\frac{C_{\mathrm{norm,int}}}{2s}
\le\frac{C_{\mathrm{norm,int}}}{s},
\]
this proves
\[
\frac{\mathsf J_{2s}^{(4,4)}(0)^2}{H_{2s}^{4,4}}
\ge c_{\mathrm{spec,int}},
\qquad
c_{\mathrm{spec,int}}
:=\frac2{3\pi C_{\mathrm{norm,int}}}>0
\qquad(s\ge1).
\]
% lean: CausalSmith.Stat.NoisydoseWeakdesignTransition.jacobi_four_even_spectral_lower

\item \emph{Interior normalization, including the smallest degrees.}
Define
\[
N_{\mathrm{even}}(m):=\max\{1,\lfloor m/8\rfloor\},
\qquad
\nu_{\mathrm{int}}(m,i)
:=2\bigl(\lceil5m/8\rceil+i\bigr),
\qquad
0\le i<N_{\mathrm{even}}(m).
\]
These are distinct positive even degrees satisfying
\[
\frac54\le\frac{\nu_{\mathrm{int}}(m,i)}m\le\frac74,
\qquad
\nu_{\mathrm{int}}(m,i)\le2m.
\]
For \(m\ge8\), the upper interval bound follows from
\[
\lceil5m/8\rceil+i
\le\frac{5m+7}{8}+\lfloor m/8\rfloor-1
\le\frac{6m-1}{8}\le\frac{7m}{8}.
\]
For \(m=4,5,6,7\), the construction selects respectively the single degree
\(6,8,8,10\), each in the stated interval.

The count satisfies
\[
N_{\mathrm{even}}(m)\ge\frac m{15}\qquad(m\ge4).
\]
Indeed, for \(m\ge8\),
\[
m\le8\lfloor m/8\rfloor+7
\le15\lfloor m/8\rfloor,
\]
and for \(4\le m\le7\) the count is \(1\).

Every summand of \(\widetilde K_m(0)\) is nonnegative. Each selected even degree contributes at least
\(c_\eta c_{\mathrm{spec,int}}\). Hence
\[
\widetilde K_m(0)
\ge
N_{\mathrm{even}}(m)c_\eta c_{\mathrm{spec,int}}
\ge c_{\mathrm{int}}m>0,
\qquad
c_{\mathrm{int}}
:=\frac{c_\eta c_{\mathrm{spec,int}}}{15}>0.
\]
% lean: CausalSmith.Stat.NoisydoseWeakdesignTransition.packetKernel_interior_lower_of_spectral_lower

\item \emph{Interior localization and derivative.}
For \(|u|\le1\), define
\[
\mathcal B_{\mathrm{int}}(m,u)
:=
\sqrt{\mathcal W_{4,4}(m;0)\mathcal W_{4,4}(m;u)}.
\]
The regularized weights satisfy
\[
\mathcal W_{4,4}(m;0)=(1+m^{-2})^9\ge1
\]
and
\[
\begin{split}
\mathcal W_{4,4}(m;u)
&=\bigl((1-u+m^{-2})(1+u+m^{-2})\bigr)^{9/2}\\
&\ge\Delta(u)^{9/2}\ge\Delta(u)^6.
\end{split}
\]
Thus
\[
\mathcal B_{\mathrm{int}}(m,u)\ge\Delta(u)^3\ge\Delta(u)^4.
\]
The angular distance dominates the coordinate distance:
\[
|u|
=
|\cos(\arccos u)-\cos(\arccos0)|
\le|\arccos u-\arccos0|
=
\delta_{\mathrm{ang}}(0,u).
\]
Localization with exponent \(2\) therefore gives
\[
\begin{split}
|\Delta(u)^4\widetilde K_m(u)|
&\le
\frac{\Delta(u)^4C_{\mathrm{loc},2}^{4,4}m}
{\mathcal B_{\mathrm{int}}(m,u)
 (1+m\delta_{\mathrm{ang}}(0,u))^2}\\
&\le C_{\mathrm{loc},2}^{4,4}m(1+m|u|)^{-2}.
\end{split}
\]
Dividing by the normalization lower bound yields
\[
|\psi_m(u)|
\le
\frac{C_{\mathrm{loc},2}^{4,4}}{c_{\mathrm{int}}}
(1+m|u|)^{-2}.
\]
This also holds outside \([-1,1]\), where the packet vanishes.

For \(|u|<1\), symmetry of the kernel and the local angular estimate, with reference and auxiliary arguments both equal to \(u\) and fixed second argument \(0\), give, for \(v\) sufficiently close to \(u\),
\[
|\widetilde K_m(v)-\widetilde K_m(u)|
\le
\frac{C_{\mathrm{lip},1}^{4,4}m^2}
{\mathcal B_{\mathrm{int}}(m,u)
 (1+m\delta_{\mathrm{ang}}(0,u))}
|\arccos v-\arccos u|.
\]
Passing to difference-quotient limits and using
\[
\left|\frac{d}{du}\arccos u\right|
=\frac1{\sqrt{\Delta(u)}}
\]
gives
\[
|\widetilde K_m'(u)|
\le
\frac{C_{\mathrm{lip},1}^{4,4}m^2}
{\mathcal B_{\mathrm{int}}(m,u)\sqrt{\Delta(u)}}.
\]
Localization with exponent \(1\) also gives
\[
|\widetilde K_m(u)|
\le
\frac{C_{\mathrm{loc},1}^{4,4}m}
{\mathcal B_{\mathrm{int}}(m,u)}.
\]
The product rule is
\[
\psi_m'(u)
=
\frac{-8u\Delta(u)^3\widetilde K_m(u)
+\Delta(u)^4\widetilde K_m'(u)}
{\widetilde K_m(0)}.
\]
Using
\[
\frac{\Delta(u)^3}{\mathcal B_{\mathrm{int}}(m,u)}\le1,
\qquad
\frac{\Delta(u)^4}{\sqrt{\Delta(u)}}\le\Delta(u)^3,
\]
we obtain
\[
\begin{split}
|\psi_m'(u)|
&\le
\frac{8C_{\mathrm{loc},1}^{4,4}m
      +C_{\mathrm{lip},1}^{4,4}m^2}
     {\widetilde K_m(0)}\\
&\le
\frac{8C_{\mathrm{loc},1}^{4,4}
      +C_{\mathrm{lip},1}^{4,4}}{c_{\mathrm{int}}}\,m.
\end{split}
\]
Here \(m\ge1\) absorbs the term independent of \(m\) after division.

On the whole line,
\[
\psi_m(u)
=
\mathcal T(u)\frac{\widetilde K_m(u)}{\widetilde K_m(0)}.
\]
The previously established regularity of \(\mathcal T\) proves that this packet is continuously differentiable. Its derivative vanishes outside \([-1,1]\), and continuity extends the bound to the two endpoints.

Set
\[
C_{\mathrm{env,int}}
:=
\frac{C_{\mathrm{loc},2}^{4,4}
      +8C_{\mathrm{loc},1}^{4,4}
      +C_{\mathrm{lip},1}^{4,4}}{c_{\mathrm{int}}}>0.
\]
We have proved, for every real \(u\),
\[
|\psi_m(u)|\le C_{\mathrm{env,int}}(1+m|u|)^{-2},
\qquad
|\psi_m'(u)|\le C_{\mathrm{env,int}}m.
\]
% lean: CausalSmith.Stat.NoisydoseWeakdesignTransition.packetinterior_bounds_of_normalization

\item \emph{Interior symmetry, masses, and normalization at zero.}
Reindexing the binomial representation in \cref{obj:lem:jacobi-binomial} by \(i\mapsto j-i\) gives
\[
\mathsf J_j^{(4,4)}(-u)
=(-1)^j\mathsf J_j^{(4,4)}(u).
\]
At \(u=0\), every odd-degree center value is therefore zero. The surviving even-degree terms give
\[
\widetilde K_m(-u)=\widetilde K_m(u),
\qquad
\psi_m(-u)=\psi_m(u).
\]

For the exponent \(\kappa=0\), we use \(|u|^0=1\), including at \(u=0\). Define
\[
C_{\mathrm{mid,int}}:=2C_{\mathrm{env,int}}>0.
\]
The envelope bounds the supremum by \(C_{\mathrm{mid,int}}\), and the derivative is bounded by \(C_{\mathrm{mid,int}}m\). Evenness and the explicit quadratic-envelope integral give
\[
\begin{split}
\int_{-1}^1|\psi_m(u)|\,du
&=
2\int_0^1|\psi_m(u)|\,du\\
&\le
2C_{\mathrm{env,int}}\int_0^1(1+mu)^{-2}\,du
\le C_{\mathrm{mid,int}}m^{-1}.
\end{split}
\]
The pointwise inequality
\(\psi_m(u)^2\le C_{\mathrm{mid,int}}|\psi_m(u)|\) then gives
\[
\int_{-1}^1\psi_m(u)^2\,du
\le
C_{\mathrm{mid,int}}\int_{-1}^1|\psi_m(u)|\,du
\le C_{\mathrm{mid,int}}^2m^{-1}.
\]
Thus, with
\[
C_{\mathrm{fin,int}}
:=
\max\{C_{\mathrm{mid,int}},C_{\mathrm{mid,int}}^2\}>0,
\]
all four analytic bounds hold with one constant. The definition and positive normalizer also give
\[
\psi_m(u)=0\quad(|u|>1),
\qquad
\psi_m(0)=\frac{\widetilde K_m(0)}{\widetilde K_m(0)}=1.
\]
% lean: CausalSmith.Stat.NoisydoseWeakdesignTransition.packetinterior_bounds

\item \emph{Interior moments and conclusion.}
Let \(j\) be an integer with \(0\le j<m\); in particular \(j\le m\). On \([-1,1]\), the taper equals the symmetric Jacobi weight:
\[
(1-u^2)^4=(1-u)^4(1+u)^4.
\]
Expanding the finite kernel sum therefore gives
\[
\begin{split}
\int_{-1}^1u^j\psi_m(u)\,du
&=
\frac1{\widetilde K_m(0)}
\sum_{\nu=0}^{2m}
\frac{\eta(\nu/m)\mathsf J_\nu^{(4,4)}(0)}
     {H_\nu^{4,4}}\\
&\qquad\qquad{}\times
\int_{-1}^1
(1-u)^4(1+u)^4
\mathsf J_\nu^{(4,4)}(u)u^j\,du
=0.
\end{split}
\]
Indeed, an inactive coefficient is zero, while an active degree satisfies
\(\nu>m\ge j\), so its integral vanishes by Jacobi orthogonality. Finite summation and compact-interval integrability justify the expansion.

The interior branch consequently satisfies every required assertion with
\[
C:=C_{\mathrm{fin,int}},
\qquad c:=1.
\]
Together with the positive-exponent branch, this proves the result for every fixed \(\kappa\in[0,2]\).
% lean: CausalSmith.Stat.NoisydoseWeakdesignTransition.weighted_packet_construction
\end{enumerate}
\end{proof}
